\begin{helper}
Let \(s,m\) be integers with \(s\ge4\) and \(m\ge1\).  Put
\[
  t=s-2,\qquad q=2^m,\qquad
  n=\frac{q^{t+1}-1}{q-1},\qquad
  d_G=\frac{q^t-1}{q-1},
\]
\[
  d_F=n-d_G,\qquad
  a=\frac{q^{t-1}-1}{q-1},\qquad
  \lambda=\sqrt{d_G-a},\qquad Q=q^{t-1},
\]
and
\[
  \eta=\max\left\{\frac{\lambda^2}{d_G^2},
    \frac{\lambda\lambda}{d_Fd_G}\right\}.
\]
Then
\[
  q^t\le n\le2q^t,\qquad
  q^{t-1}\le d_G\le2q^{t-1},\qquad
  d_F=q^t,\qquad \lambda^2=Q,
\]
\[
  \frac n2\le d_F,\qquad
  \frac{1}{16Q}\le\eta\le\frac{16}{Q}.
\]
\end{helper}
\begin{proof}
The node \(\noderef{MainTheoremPolarityParameterBounds}\) gives
\[
  q^t\le n,\qquad q^{t-1}\le d_G,\qquad d_F=q^t,\qquad
  \lambda^2=Q,
\]
and also \(\eta\le 1/Q\), because \(\eta\) is the maximum of the two displayed
spectral ratios and that maximum is bounded there by \(1/Q\).  Since \(m\ge1\),
we have \(q=2^m\ge2\).

We now sharpen the two geometric-series upper bounds.  For every integer
\(r\ge0\),
\[
  \sum_{i=0}^r q^i\le 2q^r.
\]
This is proved by induction on \(r\).  The case \(r=0\) is \(1\le2\).  If it
holds for \(r\), then
\[
  \sum_{i=0}^{r+1}q^i
    \le 2q^r+q^{r+1}
    =(2+q)q^r
    \le 2q\,q^r
    =2q^{r+1},
\]
where \(2+q\le2q\) follows from \(q\ge2\).  The standard finite geometric-sum
identity gives
\[
  n=\sum_{i=0}^t q^i,\qquad d_G=\sum_{i=0}^{t-1}q^i.
\]
Applying the induction estimate with \(r=t\) and \(r=t-1\), respectively,
gives
\[
  n\le2q^t,\qquad d_G\le2q^{t-1}.
\]
Since \(d_F=q^t\), the inequality \(n\le2q^t\) also gives \(n/2\le d_F\).

It remains to prove the lower bound for \(\eta\).  The first term in the
maximum defining \(\eta\) is at most \(\eta\), and by \(\lambda^2=Q\) it is
\[
  \frac{Q}{d_G^2}.
\]
The upper bound \(d_G\le2Q\), with \(d_G,Q>0\), implies
\[
  d_G^2\le (2Q)^2=4Q^2.
\]
Therefore
\[
  \frac{1}{4Q}\le \frac{Q}{d_G^2}\le \eta.
\]
Since \(1/(16Q)\le1/(4Q)\), we get \(1/(16Q)\le\eta\).  Finally,
\(\eta\le1/Q\le16/Q\), so the claimed two-sided eta estimate follows.
\end{proof}
